\subsection{Feasibility Before Response Alignment}
\label{sec:finite-feasibility}
\textbf{实验要求补上的逻辑层。} P1 的 matched 条件中，校准模型几乎逐轨迹
复现了失败，但所选有限映射仍只有 $3/30$ 成功。因此，校准误差、控制/映射
表达误差和目标选择偏置必须分开。小的响应预测误差可以准确预测一个坏策略；
它既不保证候选类存在可行策略，也不保证 source anchoring 优于 task-only。
以下结论条件于固定接口、已声明模型与候选类，作为方法设计与审计工具。

\textbf{可行支持与拒绝。} 设 $\Phi_F$ 是冻结的有限映射族，
$S_m\subseteq\Phi_F$ 是模式 $m$ 下满足所声明任务、约束及模式对应的映射集合。
必须区分真实的 $S_m$ 与仅在有限模型训练/验证 reset 上观察到的
$\widehat S_m$；后者不是真实安全集合的证书。
\begin{proposition}[共同映射、模式条件映射与重加权的支持边界]
对于具有正源概率的模式集合 $\mathcal M_s$，一个共同映射能保留全部模式
当且仅当 $\bigcap_{m\in\mathcal M_s}S_m\ne\varnothing$。
允许逐模式选择同一候选族中的映射时，全部模式可执行当且仅当每个
$S_m\ne\varnothing$。仅对固定映射的源模式重加权不能使其不可执行模式
变为可执行；若全部模式均无支持，则输出停止/拒绝质量，而不是归一化零权重。
\end{proposition}
\begin{proof}
共同映射必须同时属于每个集合。模式条件映射从每个非空集合选择一个元素，
反向必要性直接成立。重加权只改变固定响应支持上的概率，不产生新的命令或
闭环响应支持，故不能补回缺失的可执行模式。
\end{proof}
这是\emph{候选类内}的结论；搜索失败不证明机器人全局不可行。
P1.7 的 task-only 也允许相同逐模式选择，因此增加条件化容量本身不能
作为 response anchoring 的独立优势。有限 reset 上的全成功过滤仅是经验
拒绝规则；所有拒绝与未执行任务保留在总体分母中。

\begin{proposition}[不变执行前缀的失败上界]
设一族 transport 在到达可观测阶段边界 $\sigma$ 前，采用相同指令、
相同控制器及反馈规则。将初态、源样本和物理随机流共同耦合，则它们在
$\sigma$ 前具有相同的增广闭环轨迹。若共同前缀在到达 $\sigma$ 前
以概率 $p_{\rm pre}$ 触发不可撤销终止失败，则该族任一策略的整体成功率
至多为 $1-p_{\rm pre}$。
\end{proposition}
\begin{proof}
相同初态和随机流下，按闭环递推归纳，相同反馈指令给出相同状态、
controller carry 与事件。故前缀失败事件对全族相同；整体成功事件必须
包含于其补集。对该包含关系取概率即得。
\end{proof}
这里的 $p_{\rm pre}$ 是声明初态分布下的概率，有限模型样本只给出其经验
描述。P1.6 模型训练中出现阶段 $1$ 的回撤超力，而阶段映射子族在阶段 $2/3$
才激活，说明应检查恢复前缀的表达支持。该结论不授权放宽力限或隐去失败；
新的恢复阶段参数族须单列版本，并对所有同类控制对照开放。
另做一个不回流搜索的配对模型诊断：同一已观察的失败 reset 中，六个只改
终段的候选在 $1.04\,\mathrm s$ 全部超力，失败前位置与指令逐元素相同；
加入回撤阻抗的 $18$ 个声明候选均成功。该单 reset 模型结果检验前缀
不变条件及表达缺口，不证明跨 reset、实际目标或总体的因果收益。

\textbf{保留模式质量与停止的代价。} 设 $S$ 是给定可行性规则允许保留
的源模式集合，$p_s(S)>0$。若目标 law $q$ 只在 $S$ 上分配质量，则
\[
 \operatorname{TV}(q,p_s)\ge p_s(S^c),\qquad
 \min_{q:q(S)=1}\operatorname{KL}(q\Vert p_s)=-\log p_s(S).
\]
第二个最小值由 $q=p_s(\cdot\mid S)$ 达到：对任意支持于 $S$ 的 $q$，
$\operatorname{KL}(q\Vert p_s)=
\operatorname{KL}(q\Vert p_s(\cdot\mid S))-\log p_s(S)$。
第一个式子直接在集合 $S^c$ 上使用 TV 定义。若保留原源采样而明确对
$S^c$ 的质量停止，则该停止策略的成功率至多为 $p_s(S)$。
这些是支持限制的基本恒等式，不作为新的 OT 理论；实验中仅有模型估计
$\widehat S$，不能把它当成真实全局可行集合。特别是未通过有限训练规则
不等于该模式成功概率为零；上述停止上界描述算法选择，而不是机器人能力上界。

因此，部分能力条件下不能同时要求源 law 不变、全部模式对应和高任务成功。
P1.7 的 reweight 对照单列检验这种取舍，不能因原始 source draw 的对应下降
就判为错误。若共同优化模式 law 与映射，应把 $(q,\phi)$ 作为联合候选；
下述固定 law/同集合的任务预算保证不能自动转移到不同重加权策略。

\textbf{用任务非劣预算组织响应目标。} 固定一个非空的共同候选集合
$\mathcal A$，令真实任务与响应成本分别为 $C,R$，模型成本为
$\widehat C,\widehat R$。先求 $\widehat C_*=\min_{\phi\in\mathcal A}
\widehat C(\phi)$，再选择
\begin{equation}
 \widehat\phi_\delta\in\arg\min_{\phi\in\mathcal A}\widehat R(\phi)
 \quad\text{s.t.}\quad
 \widehat C(\phi)\le\widehat C_*+\delta.
 \label{eq:task-band-transport}
\end{equation}
任务约束、mode 条件、候选容量和模拟缓冲对完整方法与 task-only 相同。
$\delta$ 在封存目标评估前固定；它公开允许的任务成本代价，而不是通过
反复调节响应权重寻找最终测试优势。

\begin{proposition}[同一有限集合内的任务预算与响应比较]
若在 $\mathcal A$ 上 $|C-\widehat C|\le\epsilon_C$、
$|R-\widehat R|\le\epsilon_R$ 统一成立，且 Eq.~\eqref{eq:task-band-transport}
求解误差至多为 $\epsilon_{\rm opt}$（约束仍被满足），记
$C_*=\min_{\phi\in\mathcal A}C(\phi)$，则
\begin{equation}
 C(\widehat\phi_\delta)\le C_*+\delta+2\epsilon_C.
\end{equation}
若 $\delta\ge2\epsilon_C$，还有
\begin{equation}
 R(\widehat\phi_\delta)\le
 \min_{\substack{\phi\in\mathcal A\\C(\phi)\le C_*+\delta-2\epsilon_C}}
 R(\phi)+2\epsilon_R+\epsilon_{\rm opt}.
\end{equation}
\end{proposition}
\begin{proof}
以真实 task minimizer 代入模型最小值，得到
$C(\widehat\phi_\delta)\le\widehat C_*+\delta+\epsilon_C
\le C_*+\delta+2\epsilon_C$。任何满足第二式真实成本约束的候选均满足
$\widehat C(\phi)\le C_*+\delta-\epsilon_C
\le\widehat C_*+\delta$，因而属于模型预算带；再对响应成本两端应用统一
误差界和优化误差即得。
\end{proof}
该命题只比较同一候选集合。本文开发实验中的 $\delta=0.02$ 是预先固定
的经验选择规则，并未认证统一 $\epsilon_C,\epsilon_R$，所以不能把实测
小误差直接代入，宣称总体非劣或真实安全。

\textbf{模型偏差与 reset 采样误差分开。} 上式的 $\epsilon_C$ 若将有限
训练均值当作 $\widehat C$，必须同时覆盖模型偏差与 Monte Carlo 误差。
对事先固定的 $M$ 个候选和校准后独立的 $n$ 个 reset clusters，
若每 cluster 的模式平均成本属于 $[0,1]$，则以概率至少 $1-\alpha$，可以取
\[
 \epsilon_C=\epsilon_{\rm model}
       +\sqrt{\frac{\log(2M/\alpha)}{2n}}.
\]
此处 $\epsilon_{\rm model}$ 仍需外部论证；有限 occupancy 的实测误差
不是它的统一证书。即使同一 twin 精确预测实际轨迹，四个训练 reset 上
的最优候选也可能在新模型验证 reset 失败。当前 $n=4,M=24$ 的保守界
远大于 $\delta=0.02$，所以预算带只是经验设计，不能宣称已认证非劣。
在 v2 的一个新 law reset 上，full 比 task-only 多一次失败；对应重复的
实际样本均值任务成本差为 $0.06788$，尽管响应成本差为 $-0.00309$。
这不是条件命题的反例，而是缺少统一误差与充分 reset 覆盖的实测提醒。
对有限样本的精确枚举最优，也不等于对期望目标的 $\epsilon_{\rm opt}=0$。

\textbf{物理后端误差单独计入。} 设 $C_{\rm CPU}$ 为冻结目标参考环境的
期望成本，$C_{\rm twin,CPU}$ 为校准参数在同一 CPU 物理下的成本，
$C_{\rm twin,GPU}$ 为这些参数在 GPU 物理下的成本。若分别存在同一有限
候选集合上的统一界，则由三角不等式，GPU 训练均值所需的统一误差可取
\[
 \epsilon_C = \epsilon_{\rm calibration}
                  +\epsilon_{\rm engine}
                  +\epsilon_{\rm sampling},\qquad
 \epsilon_{\rm engine}\ge
 \sup_{\phi\in\mathcal A}
 |C_{\rm twin,CPU}(\phi)-C_{\rm twin,GPU}(\phi)|.
\]
这里各项是需要论证的上界，而非由少数配对轨迹自动取得的证书。
坐标系、力矩平移、状态 carry 和随机流一致只是实现层前提；接触解算差异
仍可能通过阶段切换、首次保持终止或不可恢复失败放大。双精度本身也不保证
长时闭环等价。因此 GPU 模型搜索、独立 CPU 模型验证与冻结 CPU 目标评估
应分别记账。后两者的有限验证不能被解释为对所有候选与未来 reset 的统一界。

\textbf{表达误差不能被模型精度消除。} 对更宽的参考策略类
$\Psi\supseteq\mathcal A$，定义
$\Delta_{\mathcal A}=\inf_{\phi\in\mathcal A}C(\phi)-
\inf_{\psi\in\Psi}C(\psi)\ge0$。若原标量化目标
$\widehat C+\lambda\widehat R$（$\lambda\ge0$）的优化误差为 $\epsilon_{\rm opt}$，
$R\ge0$，则以同集合真实 task minimizer $\phi_C$ 为比较点有
\begin{equation}
 C(\widehat\phi)-\inf_{\psi\in\Psi}C(\psi)
 \le\Delta_{\mathcal A}+2\epsilon_C+2\lambda\epsilon_R
   +\lambda R(\phi_C)+\epsilon_{\rm opt}.
\end{equation}
证明由模型标量化最优性、成本两端误差以及
$-\lambda R(\widehat\phi)\le0$ 直接得到。即使校准和优化完全准确，
候选表达误差与 anchoring 偏置仍可能存在；上一节的
$2\varepsilon_F+\varepsilon_{\rm opt}$ 并未消除这些项。

\textbf{接触事件裕度与策略质量分开。} 上节事件界控制真实与预测的
事件差异，不保证事件本身成功。应同时报告任务失败概率、逐约束轨迹裕度、
以及预测事件误差。P1 的三个 matched 成功在最后保持窗口中距法向力下界
仅约 $0.0008$--$0.0144\,\mathrm N$，因此均值误差、最终力或模式标签
都不足以表示执行鲁棒性。事件距离必须包含接触 witness、计数器与控制 carry；
末端位置或平均力的 RMSE 不能替代它。
同时，控制器在首次满足保持条件时终止，会使最后保持窗口的起点接近阈值。
因此上述窗口裕度只是轨迹描述，不能直接当作闭环扰动鲁棒半径；验证后者
还需要控制扰动、完整停止时刻与必要的终止后承载诊断。

\textbf{验证拒绝属于算法成本。} 令 $a_{j,m}\in\{0,1\}$ 为目标评估前冻结
的模式接受标记，$C_{j,m}\in[0,1]$ 为该映射执行时的真实任务成本。
以拒绝成本 $1$ 计入全部请求时，实际算法成本是
\[
 C_j^{\rm alg}=\sum_m p_m\{a_{j,m}C_{j,m}+1-a_{j,m}\}.
\]
因此同一源 law 下，完整方法 $F$ 与 task-only $T$ 满足
\[
 C_F^{\rm alg}-C_T^{\rm alg}
 \le \sum_{m:a_{F,m}=a_{T,m}=1}p_m(C_{F,m}-C_{T,m})
       +\sum_{m:a_{F,m}=0,\,a_{T,m}=1}p_m.
\]
共同拒绝模式贡献为零；仅 task-only 被拒绝的项非正；仅完整方法被拒绝
的项为 $p_m(1-C_{T,m})\le p_m$，即得该式。训练预算带即使控制了
映射本身的成本，也不能忽略两方法的验证拒绝差异。若同时采用不同的
模式权重 $p_F,p_T$，按 $p_F$ 应用上述比较后还需加上
$\operatorname{TV}(p_F,p_T)$，因为停止后的成本仍属于 $[0,1]$。
这区分了被执行映射的质量、接受集合和模式质量分配，防止只在成功部署
子集上宣称总体非劣。若两方法具有相同冻结映射、接受标记与源 law，则
其算法成本由结构直接相等；这不提供独立任务收益。

\textbf{边缘一致不等于源条件对应。} 令 $M\sim\mathrm{Bernoulli}(1/2)$。
响应 $Y=M$ 与 $\widetilde Y=1-M$ 的边缘 law 完全相同，但
$\Pr(Y\ne M)=0$、$\Pr(\widetilde Y\ne M)=1$。
两种策略可以具有相同成功率与目标模式频率，却具有相反的联合对应。
固定打乱 coupling 的实验验证这一度量边界；它由构造产生，不能单独证明
本文方法优于使用同一 source prior 的 task-only 控制器。

\textbf{有限类选择持平的解释。} 对冻结成本 $(C_i,R_i)$，候选 $i$
最小化 $C+\lambda R$ 当且仅当
$(C_i-C_j)+\lambda(R_i-R_j)\le0$ 对全部 $j$ 成立。
这些半直线的交集给出可直接计算的 $\lambda$ 区间，包含空区间和无界区间。
同一个映射在共同随机流下产生相同实际 law；给它不同方法名称不会产生
独立收益。该代数诊断说明何时 anchoring 不改变选择，既不是统计证书，也
不支持通过试遍 $\lambda$ 后挑选测试集胜者。

\textbf{相同控制类与额外随机性的边界。} 本文 full 与强 task-only 使用
相同参数化与信息，因此可表达相同闭环 law；source anchoring 改变的是
选择准则，不扩大相对该强对照的可执行集合。相对 reweight-only 的支持
改变不能转称相对同类控制器的表达优势。

固定一个源样本、可观测条件和有限合法动作集合，令 $c(a)$ 为已对物理
随机性取期望的标量成本。任意条件随机 kernel $T$ 满足
$\sum_a T(a)c(a)\ge\min_a c(a)$，确定性最小动作即可达到等号。
因此在可逐条件分解的期望目标中，额外 transport noise 没有自动收益；
随机源 prior 经确定性映射仍可多模态。分布耦合约束、条件覆盖需求或
不同的鲁棒博弈可能改变这一结论，必须单独分析。当前 extra-noise 对照
只是固定小扰动，并非训练过的同容量随机 kernel，不能据其单次成功数
差异证明随机 transport 必要性。
